\documentclass[a4paper,12pt]{article}

\usepackage[utf8]{inputenc}
\usepackage[ngerman]{babel}
\usepackage{geometry}
\geometry{margin=2.5cm}

\title{Vortragsskript – Multimodaler Sprach-Assistent}
\author{Ruba Shehadeh, Godfrey Uchechukwu Atulegwu, Kilian Ketelhohn}
\date{\today}

\begin{document}

\maketitle

\section*{Einleitung}
Guten Tag, ich stelle heute meinen multimodalen Sprach-Assistenten vor. 
Dieses Projekt beschäftigt sich mit der Verarbeitung gesprochener Sprache und deren Umwandlung in strukturierte Informationen und visuelle Ergebnisse.

\section*{Projektidee}
Die Grundidee meines Projekts ist es, Sprache nicht nur als Text darzustellen, sondern sie aktiv zu analysieren und visuell greifbar zu machen. 

Dafür wird gesprochene Sprache aufgenommen, in Text umgewandelt und anschließend mithilfe von KI-Technologien weiterverarbeitet.

\section*{Ziel des Projekts}
Das Ziel ist die Kombination mehrerer KI-Bereiche in einem System: Speech-to-Text, Natural Language Processing und Bildgenerierung.

Am Ende entsteht eine Pipeline, die Sprache in strukturierte Daten und anschließend in Bilder übersetzt.

\section*{Systemarchitektur}
Das System besteht aus mehreren Modulen. Zuerst wird Sprache über ein Mikrofon aufgenommen und durch Whisper in Text umgewandelt.

Danach analysiert spaCy den Text linguistisch und erkennt Wortarten wie Verben und Nomen.

Diese Informationen werden weiterverwendet, um passende Bilder zu finden oder zu generieren, die im Streamlit Interface angezeigt werden.

\section*{Speech-to-Text}
Für die Spracherkennung wird OpenAI Whisper verwendet. 
Whisper wandelt gesprochene Sprache in Text um und bildet die Grundlage für alle weiteren Verarbeitungsschritte.

\section*{Natural Language Processing}
Im nächsten Schritt wird der Text mit spaCy analysiert.

Dabei wird der Satz in einzelne Wörter zerlegt (Tokenisierung) und grammatikalisch untersucht. Besonders wichtig sind Verben und Nomen.

\section*{Informationsextraktion}
Aus der Analyse werden gezielt relevante Informationen extrahiert.

Die Nomen dienen dabei als Schlüsselbegriffe für die spätere Bildsuche oder Bildgenerierung.

\section*{Bildgenerierung}
Die extrahierten Begriffe werden genutzt, um passende Bilder zu finden.

Dafür wird die Unsplash API verwendet, die Bilder anhand von Suchbegriffen liefert.

So entsteht eine visuelle Darstellung des gesprochenen Inhalts.

\section*{Frontend}
Die Benutzeroberfläche wurde mit Streamlit umgesetzt.

Der Nutzer kann Text eingeben oder Audio hochladen. Anschließend werden Analyseergebnisse, markierte Verben und Bilder angezeigt.

\section*{Technologien}
Verwendet wurden Python, Whisper für Spracherkennung, spaCy für NLP, Streamlit für die Oberfläche und die Unsplash API.

\section*{Ergebnis}
Das Ergebnis ist eine KI-Pipeline, die Sprache in strukturierte Informationen und Bilder umwandelt.

Sprache wird damit nicht nur verstanden, sondern auch visuell interpretiert.

\section*{Fazit}
Das Projekt zeigt, wie verschiedene KI-Technologien kombiniert werden können, um ein interaktives System zu erstellen.

Es verbindet Sprachverarbeitung, Textanalyse und Bildgenerierung in einer Anwendung.

\section*{Ende}
Vielen Dank für die Aufmerksamkeit. Ich beantworte jetzt gerne Fragen.

\end{document}